El sistema visual se construye en dos niveles (ADR-9). El primero son los
\textit{tokens}: variables \acs{css} con todos los valores de color, tipografía, tamaño,
radio y sombra (\texttt{src/styles/tokens.css}). El segundo son las primitivas de
\texttt{src/ui}, componentes que usan solo esos \textit{tokens}. Los módulos funcionales
componen pantallas con las primitivas y, si necesitan estilos propios, se limitan a la
maquetación con clases prefijadas por su identificador; no pueden escribir colores ni
tipografías literales ni redefinir las clases de las primitivas. Estas reglas no dependen
de la buena voluntad de quien programa: un auditor determinista las comprueba antes de cada
\textit{commit} (\cref{SEC:TEST-AUDITS}).

\subsection{Una primera dirección rechazada}

La primera versión del sistema (v1) se inspiró en una referencia visual aportada por el
autor: fondo oscuro, tipografía monoespaciada, etiquetas en mayúsculas y un acento dorado.
Al revisarla, el autor la rechazó porque transmitía una herramienta financiera y no una
aplicación educativa. El análisis posterior mostró que la referencia se había interpretado
de forma demasiado literal: de ella solo interesaba el minimalismo (jerarquía construida con
espacio, líneas finas, pocos elementos), no su paleta ni su tipografía. Este episodio es
uno de los ejemplos más claros de decisión que no podía inferirse del código y requería a
una persona (\cref{SEC:GEN-CLASSIFICATION}).

\pendiente{Incluir capturas comparativas de la v1 rechazada y de la v2.}

\subsection{La dirección actual: clara y educativa}

La versión vigente (v2) busca la calma de un entorno de estudio, en el papel que cumple
Moodle en la universidad pero con más contención:
\begin{itemize}
  \item \textbf{Color con significado.} Fondo de papel frío (\texttt{\#f6f7f9}) y
        superficies blancas; un único color de acción, azul tinta (\texttt{\#2d4fa3}),
        para botones, enlaces, el paso actual y el foco. El rojo de corregir
        (\texttt{\#b83a2e}) se reserva para suspensos, errores y acciones destructivas; el
        verde (\texttt{\#2e7d4f}), para aprobados y tareas terminadas; el ámbar
        (\texttt{\#9a6410}), para lo pendiente o lo que requiere atención; y un gris
        azulado, para los estados de procesamiento. Cada estado de la instancia tiene
        asignado uno de estos tonos, de modo que el color de una insignia siempre significa
        lo mismo en toda la aplicación.
  \item \textbf{Tipografía para leer.} Atkinson Hyperlegible \cite{atkinson_hyperlegible},
        diseñada por el Braille Institute para maximizar la distinción entre caracteres,
        para todo el texto de la interfaz; Source Serif 4 para títulos de página y cifras
        importantes (la nota, los totales), y una monoespaciada solo para códigos
        (asignaturas, identificadores). Texto base de 15 píxeles, frases en minúscula
        inicial, sin mayúsculas sostenidas ni espaciado amplio.
  \item \textbf{Formas y profundidad.} Radios moderados para paneles y controles, bordes de
        un píxel y una única sombra, reservada a la hoja del examen, los diálogos y los
        avisos: lo único que ``flota'' es lo que de verdad está por encima.
  \item \textbf{Movimiento.} Solo como respuesta a una acción del usuario, y se respeta la
        preferencia del sistema de reducir el movimiento.
\end{itemize}

\subsection{Accesibilidad}

Se fijó un mínimo común, apoyado en las primitivas para que ninguna pantalla pueda
saltárselo: anillo de foco azul de dos píxeles en todos los controles, etiqueta asociada a
cada campo, navegación por teclado en las filas de las tablas, un color de texto secundario
mínimo para cualquier texto real y diálogos y avisos con los roles y atributos
\acs{aria} correspondientes. Se toman como referencia las pautas \acs{wcag} 2.2
\cite{wcag22}, aunque no se ha realizado una auditoría formal de conformidad
(\cref{SEC:CONC-FUTURE}).

\subsection{Diseño adaptable}

A partir de 900 píxeles de ancho, la aplicación muestra la barra superior, la navegación
lateral y el contenido; por debajo, la navegación se convierte en una franja horizontal.
Las tablas se desplazan dentro de su marco y ningún otro elemento puede desbordar la
pantalla.

\subsection{Evidencia de la arquitectura}

El cambio completo de v1 a v2 se hizo tocando únicamente los \textit{tokens}, las
primitivas y el \acs{css} del armazón. Las pantallas que Codex desarrollaba en paralelo lo
heredaron sin ninguna modificación, porque su \acs{css} propio se limitaba a unas pocas
reglas de maquetación. Para no romper las ramas en curso, los nombres de los
\textit{tokens} de v1 se mantuvieron como alias de los nuevos. La guía completa se recoge en
el \cref{AP:STYLE}.
